\section{What earlier studies changed and measured}

\label{sec:related-work}

\subsection{Which tests the repair system receives}

Smith et al. evaluated generated repairs with tests withheld during repair. A patch could satisfy the tests used to produce it and fail an independent assessment. Kong et al. varied test suites, including the number of failing tests, when measuring repair effectiveness. Motwani et al. compared choices of failing tests within a Java defect and evaluated patches using independent tests. Their design is especially relevant here because it holds the defect fixed. Our question also holds fixed the number of supplied examples and asks which of two illustrated behaviours changes under that choice. Answering it requires a separate result for each behaviour on patches from either assignment.~\cite{smith2015overfitting,kong2018impacts,motwani2022quality}

Research on repository-level agents gives other reasons to describe exactly which tests an agent sees. Ahmed et al. investigated test overfitting on SWE-bench. Their manuscript includes an experimental condition in which the agent could see developer tests ordinarily held back. TDFlow studied workflows with human-written and system-generated tests. Access and authorship can each change what a test contributes to a repair. In the proposed comparison, each arm receives one example drawn from the same audited finding. A difference between arms can then be expressed in terms of the two behaviours the examples illustrate.~\cite{ahmed2026overfitting,han2026tdflow}

\subsection{What counts as a corrected behaviour}

Security repair benchmarks offer several ways to define a successful patch. Vul4J pairs historical Java revisions with tests that help establish reproducible cases. APR4Vul used manual assessment to examine both vulnerability removal and side effects on Java functionality. SEC-bench includes patching among the agent tasks it evaluates on real security problems. EVMbench is the closest benchmark in the Solidity setting. Its patch task already assesses existing project tests and additional evidence withheld from the agent. Hidden assessments and project tests therefore have a direct precedent in this setting.~\cite{bui2022vul4j,bui2024apr4vul,lee2025secbench,wang2026evmbench}

The choice of what a test observes also affects a verdict. EVMPatch used historical transactions to assess intended functionality after patching. Bobadilla et al. examined this concern for automated smart-contract fixes. In their PatchBench preprint, Shen et al. compare crash-only grading with checks of security and semantics in C and C++. Ye et al. generated tests using a developer patch as the oracle. Nilizadeh et al. evaluated Java patches against formal behavioural specifications. Our historical revision illustrates a permitted correction. A patch may satisfy the recorded requirement by another implementation or by producing another explicitly permitted result. Each verdict is limited to the operations for which we can establish that requirement.~\cite{rodler2021evmpatch,bobadilla2026fixes,shen2026patchbench,ye2021assessment,nilizadeh2021overfitting}

\subsection{The unanswered comparison}

A targeted literature search completed on 23 November 2026 compared fourteen primary works on test choice, security benchmarks, and patch assessment. Among those materials, we found no study that randomized which of two behaviours of the same audited Solidity defect an example illustrated and assessed both behaviours under both assignments. Motwani et al. are close in their within-defect variation of tests, and EVMbench in its use of withheld grading. The remaining distinction has a practical consequence. If one behaviour is easier to repair, it could succeed more often in both arms. Grading it under both assignments gives a comparison that holds its difficulty fixed. For example, if a patch addresses both obligations in either arm, it contributes evidence of a useful repair without favouring either assignment. If patches generally address only the easier obligation, the same logic applies. The causal question concerns a change in each behaviour between assignments, quite apart from its overall frequency. That comparison requires evidence that the behaviours can differ.~\cite{motwani2022quality,wang2026evmbench}
